\documentclass[12pt]{article}
\usepackage[utf8]{inputenc}
\usepackage[T1]{fontenc}
\usepackage{hyperref}
\title{Pa'anci Language Resources for K-12 Curriculum Development}
\author{Compiled by Claude Sonnet 5.5 (claude-sonnet-5-5) for Andr\'as Kornai}
\date{}
\begin{document}
\maketitle
\section{Introduction}
Pa'anci (Pa'a, also Afa, Afawa, Fucaka) is a Chadic (Afro-Asiatic) language of Bauchi State, Nigeria. Speaker estimates range from about 8,000 to 34,000, and speakers are reported to be shifting to Hausa. This search found only general reference pages and no educational materials.
\section{Primary Educational Resources}
No primers, textbooks or school programmes were found in the search results.
\section{Linguistic and Reference Materials}
\begin{itemize}
\item Wikipedia article on the Pa'a language: \url{https://wikipedia.com/wiki/Paa_language}
\item Endangered Languages Project entry: \url{https://endangeredlanguages.com/lang/10445}
\end{itemize}
\section{Religious and Community Resources}
\begin{itemize}
\item Joshua Project language page (population and language-shift notes): \url{https://joshuaproject.net/index.php/languages/pqa}
\end{itemize}
\section{Digital and Audio Resources}
None were found.
\section{Curriculum Development Context}
Very little is documented online. Any curriculum would require new fieldwork and community partnership, with attention to the pressure from Hausa.
\end{document}